% 中文节标题：相关工作
\section{Related Work}

% 中文段落标题：全局人体运动重建。
% 人体网格恢复已经从参数化人体模型拟合，逐渐发展为从单目图像和视频直接进行前馈式回归。基于图像的方法主要在相机坐标系中估计人体姿态和形状，而时序模型与全画面多人方法进一步改善了运动一致性以及拥挤场景中的多人推理能力。然而，当相机本身发生运动时，相机坐标系下的预测无法直接描述人体在三维环境中的全局运动。为此，现有方法通过运动先验与序列优化、人体—相机联合建模，或透视关系与尺度标定，共同恢复人体轨迹和相机运动；近期系统还将人体、相机和周围场景几何纳入统一重建。其中，离线方法能够利用完整序列提供的全局约束，但通常需要在整段视频上反复优化人体运动、相机轨迹和场景几何，带来较高的推理延迟与计算开销，也难以扩展到不断增长的长视频。统一前馈方法降低了逐序列优化的成本；对于持续视频流，模型还需要仅根据当前与历史观测在线更新重建结果，并保持有界的历史表示规模和稳定的单步计算开销。因此，问题进一步从如何利用完整视频恢复全局关系，转向如何通过持续维护的历史表示实现在线重建。
\paragraph{Global human motion reconstruction.}
Human mesh recovery has evolved from fitting parametric body models to
feed-forward regression from monocular images and
videos~\citep{kanazawa2018hmr,kocabas2020vibe,goel2023humans4d}. Image-based
methods primarily estimate human pose and shape in the camera frame, while
temporal models and full-frame multi-person approaches improve motion
consistency and reasoning in crowded
scenes~\citep{sun2021romp,sun2022bev,sun2023trace,baradel2024multihmr,wang2025prompthmr,sun2024aios}.
When the camera itself moves, however, camera-frame predictions do not directly
describe global human motion in the surrounding 3D environment. Existing
methods therefore recover human trajectories together with camera motion
through motion priors and sequence
optimization~\citep{yuan2022glamr,ye2023slahmr,kocabas2024pace,li2024coin},
joint human--camera
modeling~\citep{wang2024tram,shin2024wham,shen2024gvhmr,yin2024whac}, or
perspective and scale
calibration~\citep{patel2025camerahmr,yang2025checkerboards}. More recent
systems further incorporate humans, cameras, and surrounding scene geometry
into a unified
reconstruction~\citep{muller2025reconstructing,liu2026josh,zhang2025odhsr,li2026unish,rojas2025hamst3r,yang2026feedforward,liu2026trophies}.
Offline methods can exploit global constraints from a complete sequence, but
typically require iterative joint optimization of human motion, camera
trajectories, and scene geometry over the entire video, incurring substantial
latency and computation and limiting scalability to growing sequences.
Unified feed-forward methods reduce the cost of per-sequence optimization;
continuous video streams additionally require reconstruction to be updated
online using only current and past observations, with a bounded history
representation and stable per-step computation. The question thus shifts from
recovering global relationships from a complete video to maintaining them
through an evolving history representation during online inference.

% 中文段落标题：在线重建。
% 围绕这一目标，在线重建按照时间顺序处理视频，并持续更新相机、几何与人体估计。传统SLAM、学习式稠密建图和前馈点图模型分别为增量式相机定位、场景建图和直接三维预测奠定了基础；近期流式方法进一步通过空间记忆、渐进式配准或递归状态保留历史观测。在人体重建方面，OnlineHMR研究世界坐标系下的在线人体恢复，Human3R则以CUT3R的递归重建为基础，将在线恢复扩展到相机、场景和多个人体。在此基础上，UniCon3R和GUSH3R分别从人体—场景接触与可渲染动态表示等方向扩展统一重建，其他流式三维方法则进一步研究长期记忆、动态几何和递归状态的稳定更新。然而，现有在线方法主要面向连续拍摄的单目视频，通常依赖相邻帧之间稳定的时间关系和视觉重叠来延续已有重建。多镜头视频中的镜头切换会突然改变相机视角与可见内容，使相邻镜头不再满足这种连续性假设，但不同镜头的重建结果仍需被组织到统一的空间关系中。因此，如何跨越镜头切换维持重建的一致性，构成了在线重建进一步面向多镜头视频时需要解决的问题。
\paragraph{Online reconstruction.}
To this end, online reconstruction processes video in temporal order and
continually updates estimates of cameras, geometry, and people. Classical
SLAM, learned dense mapping, and feed-forward point-map models provide the
foundations for incremental camera localization, scene mapping, and direct 3D
prediction, respectively~\citep{teed2021droidslam,sun2021neuralrecon,wang2024dust3r,leroy2024mast3r,zhang2025monst3r,wang2025vggt}.
Recent streaming methods further preserve past observations through spatial
memory, progressive registration, or a recurrent
state~\citep{wang2025spann3r,liu2025slam3r,wang2025cut3r}. For human
reconstruction, OnlineHMR studies online recovery in world
coordinates~\citep{zhao2026onlinehmr}, while \humanthree{} builds on CUT3R's
recurrent reconstruction to jointly recover cameras, scenes, and multiple
people online~\citep{chen2026human3r}. Building on this line, UniCon3R and
GUSH3R extend unified reconstruction through human--scene contact and
renderable dynamic representations,
respectively~\citep{sur2026unicon3r,abe2026gush3r}; other streaming 3D methods
study long-term memory, dynamic geometry, and stable recurrent
updates~\citep{chen2025long3r,wu2025point3r,zhuo2025st4r,chen2026ttt3r,zheng2026ttsa3r,li2026recal3r}.
Existing online methods, however, are primarily designed for continuously
captured monocular video and typically rely on stable temporal relations and
visual overlap between adjacent frames to carry the reconstruction forward. A
shot transition abruptly changes the camera viewpoint and visible content,
invalidating this continuity assumption even though reconstructions from
different shots must still share a coherent spatial relation. Extending online
reconstruction to multi-shot video therefore requires consistency to be
maintained across shot transitions.

% 中文段落标题：多镜头人体重建。
% 多镜头人体重建旨在从包含镜头切换的单目视频中恢复空间一致的多人体三维运动，并保持跨镜头的人物身份一致性。早期工作分别研究了影视内容中的人物跟踪与重识别、单人体跨镜头恢复以及演员—环境联合重建。HumanMM通过相机运动估计与序列级运动整合，在多个镜头之间对齐人体朝向和全局轨迹，但主要面向单人体视频。与本文任务最接近的Multi-THuMBS利用三维场景先验将镜头边界两侧的观测重建到共同空间中，并结合外观、姿态和几何线索完成人物关联，随后通过轨迹平滑与跨相机重投影约束优化全局结果；该方法将多镜头重建扩展到多人场景，但采用依赖已观测序列的多阶段全局优化。ShowMak3R同样处理多人物和多镜头内容，但通过逐视频优化人物与场景的神经表示，主要面向可渲染的影视内容重建，而非在线人体运动恢复。此外，跨镜头身份保持也与三维感知跟踪以及基于外观、姿态或关键点的人物重识别相关。这些工作从全局运动恢复、空间对齐和人物关联等角度推进了多镜头人体重建，但多采用镜头边界两侧的联合处理或对完整序列进行逐视频优化，尚未讨论在线递归重建中持续维护的历史信息应当如何跨越镜头边界。Shot3R面向这一流式设定，将历史信息暂时用于建立镜头间空间关系，但不再将其作为新镜头的递归状态持续传播，从而在顺序处理多镜头视频的同时维持空间与身份一致性。
\paragraph{Multi-shot human reconstruction.}
Multi-shot human reconstruction aims to recover spatially consistent
multi-person 3D motion from monocular videos with shot transitions while
preserving identities across shots. Early work separately studied person
tracking and re-identification in TV content, single-person recovery across
shots, and joint actor--environment
reconstruction~\citep{pavlakos2022tvtracking,pavlakos2022multishot,pavlakos2022tvreconstruction}.
HumanMM aligns human orientation and global trajectories across shots through
camera-motion estimation and sequence-level motion integration, but focuses on
single-person videos~\citep{zhang2025humanmm}. Most closely related to our
setting, Multi-THuMBS uses a 3D scene prior to reconstruct observations on both
sides of a shot boundary in a common space, associates people using
appearance, pose, and geometric cues, and then optimizes the global result with
trajectory smoothing and a cross-camera reprojection
constraint~\citep{on2026multithumbs}. It extends multi-shot reconstruction to
multiple people, but adopts multi-stage global optimization over an observed
sequence. ShowMak3R also handles multi-person, multi-shot content, but
optimizes per-video neural representations of actors and scenes, targeting
renderable cinematic reconstruction rather than online human-motion
recovery~\citep{kim2025showmak3r}. Cross-shot identity persistence is further
related to 3D-aware tracking and appearance-, pose-, or keypoint-based
re-identification~\citep{rajasegaran2022phalp,he2021transreid,yuan2025pose2id,somers2024kpr}.
Together, these works advance multi-shot human reconstruction through global
motion recovery, spatial alignment, and person association, but commonly
jointly process observations across the shot boundary or optimize a complete
sequence. They therefore leave open how a maintained history representation
should cross a shot boundary during online recurrent reconstruction. \method{}
addresses this streaming setting by using history temporarily to establish the
inter-shot spatial relation, without propagating it as the recurrent state of
the new shot, thereby maintaining spatial and identity consistency under
sequential multi-shot processing.
